\section{The Stable Reference Law, Smooth Examples and SGD}
\label{app:reference}

The following condition proves that stationary risk shrinks with the
step size for a broader class than the two-matrix example. Let the bounded positive-definite matrix $A$ be
independent of the complete bounded noise vector $e$, let
$A-\bar H$ and $\bar H-A$ have the same law, and retain coordinate sign
invariance and equal nonzero noise marginals. This central symmetry means that the deviation matrix from the mean
has the same law as its negative. It does not mean symmetry of $A$ about
zero, and it allows noncommuting matrices.

For fixed $u$, put $m_i=(\bar H u)_i$ and
$\xi_i(u)=((A-\bar H)u)_i+e_i$. Its two summands are independent symmetric variables, so their sum is
symmetric. Across time the coordinate gradients are independent with
common law $m_i+\xi_i(u)$. For one numerator summand condition on all
other times and their squared-energy term $J$. Its own-time scalar action
$x/(\sqrt{d_jx^2+J}+\epsilon)$ is odd increasing. Its mean at the
symmetric $\xi_i(u)$ is zero; shifting by $m_i$ gives its sign.
Summing with the nonnegative history weights proves
\[
 m_iF_i(u)\geq0,\qquad
 u^\top\bar H G(u)\geq\lambda u^\top\bar H u.
\]
The distribution of $\xi_i(u)$ can vary with $u$; the argument fixes $u$
before averaging its history.

The exact rowwise identity and independence give $DF(0)=c\bar H$,
$c=\E\Psi(|e_i|)>0$. Flipping the entire noise history of one coordinate eliminates every
cross-coordinate covariance at zero. Equal marginal distributions give
fluctuation covariance $DI$, with $D>0$. Lemma~\ref{lem:regular-transfer} applies with
$P_0=\bar H$, $c_0=\lambda$. The drift matrix
$B=\lambda I+c\bar H$ is symmetric positive definite, and the unique
OU covariance is $S=D B^{-1}/2$. Thus
\[
 \lim\frac{\E L(u)}\eta
 =\tfrac12\operatorname{tr}(\bar H S)
 =\tfrac D4\operatorname{tr}[\bar H(\lambda I+c\bar H)^{-1}]>0.
\]
All positive absolute scaled moments converge by the same lemma. This
proves \eqref{eq:reference-limit} and the independent half of
Theorem~\ref{thm:main}.

\paragraph{A fixed smooth joint law.}
The unsmoothed independent Hessian takes $sH_1,sH_2$ with equal probabilities,
so its deviation from $\bar H$ is a symmetric two-point matrix law. Add
$\tau B_0$ with an independent radial $C^\infty$ bump density supported
on a Frobenius ball in $\operatorname{Sym}_2$. Its mean is zero, it is
centrally symmetric, and it is invariant under coordinate exchange. Add
$\tau Z_0$ to $e$ with independent iid symmetric smooth compact coordinate
jitters. Use these same independent kernels in both copulas.

The full Hessian and noise distributions still match separately. The
joint law is unchanged by flipping any coordinate of $e$ while keeping
$A$ fixed, so the matrix/noise mixed moments still vanish. Swapping the two
coordinates in both the matrix and noise vector also preserves the law. In the independent arm, the matrix
and noise vector remain independent, and the centered matrix law
retains its central symmetry. Both joint laws are finite mixtures
of translated smooth product densities on
$\operatorname{Sym}_2\times\R^2$. Their supports are compact and have nonempty interior in all matrix
and noise coordinates. Choose one sufficiently small $\tau>0$ so that all sample
Hessians remain uniformly positive definite.

At fixed $\epsilon$, coupling these jitters to their unsmoothed marks and
using the bounded continuous Hilbert quotient derivative proves convergence
of the origin Jacobians as $\tau\downarrow0$ by dominated convergence.
Their strictly positive and negative eigenvalue margins therefore
persist for a sufficiently small fixed positive $\tau$. Coordinate-swap symmetry keeps the coupled antisymmetric direction
invariant. Its drift is zero on the symmetric line and repels near zero,
which gives the hypotheses of
Proposition~\ref{prop:matrix-stationary-repulsion}. The conditional noise bound continues to apply to the smooth
sign-invariant noise. The coupled risk has a positive lower bound; the
bounded optimizer state gives its finite upper bound. This proves the two risk orders.

For two independent matrices with these full-dimensional smooth laws,
commutation has probability zero. For a symmetric matrix with distinct
eigenvalues, the symmetric matrices that commute with it lie in a
lower-dimensional subspace. Matrices with repeated eigenvalues themselves
form a measure-zero algebraic set. All smoothing
parameters and problem laws were fixed before the step limit.

\paragraph{Exact SGD covariance equality.}
For $0<\eta\leq1/(h_{\max}+\lambda)$, set
$T=I-\eta(A+\lambda I)$ and $Q=\E ee^\top$. Every sample map shrinks distances uniformly, with norm at most $1-\eta(h_{\min}+\lambda)<1$.
Its stationary mean is zero, and freshness and $\E[A_{ij}e_k]=0$ give
the exact covariance recurrence
\[
 \Sigma=\mathcal L_\eta(\Sigma)+\eta^2Q,
 \qquad \mathcal L_\eta(S)=\E[TST^\top].
\]
The operator has norm at most the squared common contraction factor, so
$(I-\mathcal L_\eta)^{-1}$ exists and the solution is unique. Both the covariance-update operator
and the added-noise term depend only on the matched complete separate
distributions. Thus the covariance
and quadratic risk agree exactly between copulas, despite noncommutation.
Moreover
\[
 \operatorname{tr}\Sigma\leq
 \frac{\eta^2\operatorname{tr}Q}
 {1-(1-\eta(h_{\min}+\lambda))^2}=O(\eta).
\]
Expanding the recurrence and dividing by $\eta$ shows that
$\Sigma/\eta\to\Xi$, where
\[
 (\bar H+\lambda I)\Xi+\Xi(\bar H+\lambda I)=Q.
\]
The solution is unique and positive definite because $Q$ is positive
definite in this construction. Hence both matched SGD risks have order
$\eta$. This completes Theorem~\ref{thm:main}; equality of full invariant
distributions is not asserted.

\paragraph{A completely specified smooth counterexample.}
Use the atomic matrices and optimizer constants displayed in
Section~\ref{sec:matrix-experiment}. This geometry admits an explicit fixed
smoothing scale.
Write $H_{\min},H_{\max}$ for its minimum and maximum sample eigenvalues.
Let $B$ have normalized density on $\operatorname{Sym}_2$ proportional to
\[
 \exp\!\left[-\frac1{1-\|B\|_F^2}\right]
                         \ind\{\|B\|_F<1\},
\]
and let $Z_1,Z_2$ be independent with normalized density proportional to
$\exp[-1/(1-z^2)]\ind\{|z|<1\}$, independent of all original marks and
of $B$. Choose the fixed dyadic number $\tau=2^{-20}$ and replace
$A,e$ by $A_\tau=A+\tau B$, $e_\tau=e+\tau Z$ in both laws.
These smooth densities are centrally symmetric; the matrix law is also
exchange invariant. Their bounds are $\|B\|_{\rm op}\leq1$, $|Z_i|\leq1$.

To certify that this particular scale preserves the strict margins, let
$T_\epsilon$ be the Hilbert quotient in Appendix~\ref{app:transfer}.
Its derivative norm is at most $C_*/\epsilon$ and its derivative has
Lipschitz constant $6C_*/\epsilon^2$. At zero, the original coordinate
Hessian-row history map has operator norm at most $H_{\max}$; its
perturbation has norm at most $\tau$, and the noise-history perturbation
has Hilbert norm at most $\tau$. The chain rule therefore gives
\[
 \|D_uf_{\tau,i}(0)-D_uf_{0,i}(0)\|
 \leq\tau C_*\left(\frac{6H_{\max}}{\epsilon^2}
                                    +\frac1\epsilon\right).
\]
Expectation and the two row norm squares imply
\begin{equation}\label{eq:explicit-smoothing-margin}
 \|J_\tau-J_0\|_{\rm op}
 \leq\sqrt2\tau C_*\left(\frac{6H_{\max}}{\epsilon^2}
                                    +\frac1\epsilon\right)
 <\frac{99082.5}{1048576}<.1.
\end{equation}
The numerical bound uses only $H_{\max}<4$, $C_*<1.1$, $\epsilon>.02$
and $\sqrt2<1.5$. Its last inequality is the exact integer comparison
$1981650<2097152$. It is a uniform derivative bound, not a fitted or
simulation-based smoothing tolerance.

Exchange equivariance makes $J_\tau$ symmetric with the same symmetric
and antisymmetric eigendirections. Thus each decay-plus-field eigenvalue
changes by less than $.1$. The coupled antisymmetric value remains below
$-.302334$, while its symmetric value remains above $6.685694$.
The independent antisymmetric and symmetric values remain above
$2.186275$ and $9.473861$, respectively. The eigenvalue perturbation bound known as Weyl's inequality and the
atomic bounds $H_{\min}>.399$, $H_{\max}<3.601$ give
\[
 \lambda_{\min}(A_\tau)>.398999,\qquad
 \lambda_{\max}(A_\tau)<3.601001,\qquad
 \operatorname{cond}(A_\tau)<10.
\]
The population eigenvalues stay $.75,3.25$. All complete separate distributions and
mixed-moment matching identities remain unchanged. The independent arm retains
independence of its complete Hessian/noise vectors and central symmetry
of its Hessian law. The coupled arm retains conditional noise-sign symmetry, a drift that
vanishes on the symmetric line, and the uniform conditional noise bound.

Both joint laws now have specified smooth densities in all matrix and
noise coordinates,
and the restoring and repulsion arguments above establish the stationary
orders $\Theta(\eta)$ and $\Theta(1)$ for this fixed pair. The smoothing does not change
with the step size. The numerical trajectories elsewhere use the
atomic fixture; they are not mislabeled as smooth-law simulations.
The bound proves that the chosen smoothing scale and sample condition
numbers are valid. It gives no numerical value for the stationary-risk
floor, time needed to escape the origin, or convergence rate to stationarity.
